\section*{Chapitre \ref{ch:g6:water-underground} --- L'eau souterraine}

\begin{solution}{exo:g6:water-underground:1}
\omterm{def:g6:water-underground:permeability}{Perméables}~: \omterm{def:g3:stones-pebbles-sand:sand}{sable}, \omterm{def:g3:stones-pebbles-sand:pebble}{graviers}, craie. \omterm{def:g6:water-underground:permeability}{Imperméables}~: argile, granite sans
fissures, argilite.
\end{solution}

\begin{solution}{exo:g6:water-underground:2}
(a) \omterm{def:g6:water-underground:water-table}{niveau de la nappe}~; (b) \omterm{def:g6:water-underground:aquifer}{aquifère}~; (c) \omterm{def:g6:water-underground:spring}{source}~; (d) \omterm{def:g6:water-underground:spring}{puits}.
\end{solution}

\begin{solution}{exo:g6:water-underground:3}
(a) Elle s'évapore et retourne dans l'air~; elle ruisselle à la surface
jusque dans les ruisseaux~; elle s'infiltre dans le sol (\omterm{def:g6:water-underground:infiltration}{infiltration}).
(b) Sur le parking~: sa surface dure ne laisse presque pas d'eau
s'infiltrer, donc l'eau ruisselle. Dans la prairie, le sol et l'herbe en
laissent s'infiltrer une grande partie.
\end{solution}

\begin{solution}{exo:g6:water-underground:4}
(a) Faux~: l'argile est très \omterm{def:g6:water-underground:porosity}{poreuse} mais \omterm{def:g6:water-underground:permeability}{imperméable}, parce que ses \omterm{def:g6:water-underground:porosity}{pores}
sont minuscules. (b) Vrai. (c) Faux~: il suit la forme du terrain, en plus
doux, et il monte et descend avec les saisons. (d) Vrai~: environ 30~\% de
toute l'\omterm{def:g2:water-on-earth:salt-fresh}{eau douce} est de l'\omterm{def:g6:water-underground:groundwater}{eau souterraine}, contre environ un quart de
pour cent dans les \omterm{def:g1:landscapes-around-us:water}{lacs}.
\end{solution}

\begin{solution}{exo:g6:water-underground:5}
$10 \times 25 \div 100 = 2.5$~: tu peux verser 2.5~litres avant que le seau
ne déborde. $10 \times 22 \div 100 = 2.2$~: 2.2~litres s'écouleraient par
le trou.
\end{solution}

\begin{solution}{exo:g6:water-underground:6}
(a) $2 \times 1 \times 1 = 2$~mètres cubes, soit 2000~litres.
(b) $2000 \times 20 \div 100 = 400$~litres. (c) $2000 \times 18 \div 100 =
360$~litres s'écoulent~; $400 - 360 = 40$~litres restent collés.
\end{solution}

\begin{solution}{exo:g6:water-underground:7}
(a) La \omterm{def:g6:water-underground:spring}{source} est à environ 203~m, la \omterm{def:g1:landscapes-around-us:water}{rivière} à environ 132~m. (b) La
\omterm{def:g6:water-underground:spring}{source} est environ 70~m plus haut que la \omterm{def:g1:landscapes-around-us:water}{rivière}~: son eau descend la
pente jusqu'à la \omterm{def:g1:landscapes-around-us:water}{rivière}. (c) À l'intérieur du versant de craie, le \omterm{def:g6:water-underground:water-table}{niveau
de la nappe} est sous le sol~; il ne rencontre la surface que là où la base
de la craie, posée sur l'argilite \omterm{def:g6:water-underground:permeability}{imperméable}, affleure sur la pente.
\end{solution}

\begin{solution}{exo:g6:water-underground:8}
(a) Après la saison des \omterm{def:g4:clouds-rain-snow:precipitation}{pluies}~: $401 - 336 = 65$~m~; après la saison
sèche~: $401 - 314 = 87$~m. (b) $336 - 314 = 22$~m.
\end{solution}

\begin{solution}{exo:g6:water-underground:9}
Plus profond. Le \omterm{def:g6:water-underground:water-table}{niveau de la nappe} suit la forme du terrain, mais en plus
doux~: au sommet de la \omterm{def:g1:landscapes-around-us:landform}{colline}, il est plus haut que sous le \omterm{def:g6:water-underground:spring}{puits}~A, mais
le sol monte beaucoup plus, si bien que l'eau se trouve plus loin sous la
surface (sur la coupe, à plus de 110~m de profondeur).
\end{solution}

\begin{solution}{exo:g6:water-underground:10}
Les \omterm{def:g6:water-underground:porosity}{pores} de l'argile sont minuscules~: l'eau reste collée aux paillettes
et ne peut pas traverser l'argile pour aller vers le \omterm{def:g6:water-underground:spring}{puits}. L'argile est
\omterm{def:g6:water-underground:porosity}{poreuse} mais \omterm{def:g6:water-underground:permeability}{imperméable}~: seule une petite quantité d'eau (20~litres sur
500) s'en écoule.
\end{solution}

\begin{solution}{exo:g6:water-underground:11}
(a) $30.1 \div 0.26 \approx 116$~: plus de cent fois plus.
(b) $100 - 69 = 31$~: il reste environ 31~\%, dont environ 30 pour l'\omterm{def:g6:water-underground:groundwater}{eau
souterraine} à elle seule. (c) C'est de loin le plus grand stock d'\omterm{def:g2:water-on-earth:salt-fresh}{eau douce} liquide~; la
glace est gelée, et elle se trouve surtout loin, dans les régions polaires
et sur les \omterm{def:g1:landscapes-around-us:landform}{montagnes}. Et l'\omterm{def:g6:water-underground:groundwater}{eau souterraine} se trouve presque partout sous
nos pieds.
\end{solution}

\begin{solution}{exo:g6:water-underground:12}
Les deux \omterm{def:g6:water-underground:spring}{puits} puisent dans le même \omterm{def:g6:water-underground:aquifer}{aquifère} de craie. La pompe puissante
a retiré l'eau plus vite que la \omterm{def:g4:clouds-rain-snow:precipitation}{pluie} ne recharge la craie~: le \omterm{def:g6:water-underground:water-table}{niveau de
la nappe} a donc baissé tout autour, jusque sous le fond du second \omterm{def:g6:water-underground:spring}{puits},
qui s'est tari. Les agriculteurs pourraient se mettre d'accord pour pomper
moins, surtout en été, partager un même \omterm{def:g6:water-underground:spring}{puits}, ou approfondir le second
\omterm{def:g6:water-underground:spring}{puits}~; surtout, ils ne devraient pas prendre plus d'eau que la \omterm{def:g4:clouds-rain-snow:precipitation}{pluie} n'en
rapporte.
\end{solution}

\begin{solution}{pb:g6:water-underground:1}
\textbf{1.} La ferme est posée sur la craie~; dessous se trouve
l'argilite. La craie est \omterm{def:g6:water-underground:permeability}{perméable}, l'argilite ne l'est pas.

\textbf{2.} La \omterm{def:g4:clouds-rain-snow:precipitation}{pluie} s'infiltre dans la craie et s'arrête sur l'argilite.
Là où la base de la craie affleure sur la pente, le \omterm{def:g6:water-underground:water-table}{niveau de la nappe}
rencontre le sol et l'eau sort~: c'est la \omterm{def:g6:water-underground:spring}{source}.

\textbf{3.} $401 - 336 = 65$~m sous le sol.

\textbf{4.} Le fond du \omterm{def:g6:water-underground:spring}{puits} est à $401 - 100 = 301$~m, sous le \omterm{def:g6:water-underground:water-table}{niveau de la
nappe}, à 314~m~: le \omterm{def:g6:water-underground:spring}{puits} atteint encore l'eau, de $314 - 301 =
13$~m.

\textbf{5.} $100 \times 50 \times 2 = 10\,000$~mètres cubes.

\textbf{6.} $10\,000 \times 20 \div 100 = 2000$~mètres cubes.

\textbf{7.} $10\,000 \times 18 \div 100 = 1800$~mètres cubes s'écoulent~;
$2000 - 1800 = 200$~mètres cubes restent collés.

\textbf{8.} $1800 \times 1000 = 1\,800\,000$~litres.

\textbf{9.} $18 \times 1000 = 18\,000$~litres par jour.

\textbf{10.} $18 \times 7 = 126$~mètres cubes par semaine.

\textbf{11.} La \omterm{def:g4:clouds-rain-snow:precipitation}{pluie} de la saison humide s'infiltre dans la craie et la
recharge chaque année (tant que les fermes ne pompent pas plus que ce que
la \omterm{def:g4:clouds-rain-snow:precipitation}{pluie} rapporte).

\textbf{12.} $1800 \div 18 = $~\textbf{100 jours}~: un peu plus de trois
mois sans aucune \omterm{def:g4:clouds-rain-snow:precipitation}{pluie}.
\end{solution}
